\chapter{भार बनाम द्रव्यमान: बलों का मापन}\label{ch:g9:weight-vs-mass}

पाँच साल पहले, काग़ज़ पर बने \omterm{def:g2:moon-in-the-sky:moon}{चंद्रमा} पर, तुमने ``भारी'' को दो हिस्सों
में बाँटा था: सामान के लिए \omterm{def:g3:measuring-mass:mass}{द्रव्यमान}, और खिंचाव के लिए \omterm{def:g4:weight-and-mass:weight}{भार}। वह बँटवारा
समझदारी था; अब वह अंकगणित बन जाता है। बलों को उनका \omterm{def:g6:measurement-in-science:measuring}{मात्रक} मिलता है ---
गिरते सेब वाले उसी शख़्स के नाम पर --- \omterm{def:g4:weight-and-mass:weight}{भार} को उसका सूत्र, और बचपन वाली
रबर बैंड की तराज़ू पेशेवर पोशाक में लौट आती है।

\section{न्यूटन}

\begin{definition}[न्यूटन]\label{def:g9:weight-vs-mass:newton}
\omterm{def:g2:pushes-and-pulls:force}{बल} --- धक्के, खिंचाव, \omterm{def:g4:weight-and-mass:weight}{भार}, \omterm{def:g9:gravitation:gravitation}{गुरुत्वाकर्षण} की पकड़ --- \emph{न्यूटन}%
\index{न्यूटन} (\unit{N}) में मापे जाते हैं, जो \omterm{def:g9:gravitation:gravitation}{गुरुत्वाकर्षण} के नियम
के लेखक के सम्मान में है। काम के लंगर: एक न्यूटन क़रीब एक छोटे सेब के
\omterm{def:g4:weight-and-mass:weight}{भार} जितना है (इतिहास मुस्कुरा देता है); तुम्हारे बस्ते का \omterm{def:g4:weight-and-mass:weight}{भार} कोई
\qty{50}{N}; एक मज़बूत हाथ मिलाना क़रीब \qty{100}{N} से दबाता है; और
पिछले अध्याय वाला \omterm{def:g2:moon-in-the-sky:moon}{चंद्रमा} का पट्टा $10^{20}$ तक चला गया था।
\end{definition}

\begin{definition}[बलमापी]\label{def:g9:weight-vs-mass:dynamometer}
\emph{बलमापी}\index{बलमापी} \omterm{def:g2:pushes-and-pulls:force}{बल} नापने वाला यंत्र है: किसी खोल में बैठी
एक अंशांकित कमानी, जो लगाए गए \omterm{def:g2:pushes-and-pulls:force}{बल} के अनुपात में खिंचती है और जिसका
सूचक सीधे \omterm{def:g9:weight-vs-mass:newton}{न्यूटन} पढ़ा देता है। यह तुम्हारी रबर बैंड वाली तराज़ू का बड़ा
हुआ रूप है --- और अपने से पहले आए \omterm{def:g8:voltage-voltmeter:voltmeter}{वोल्टमीटर} तथा \omterm{def:g8:intensity-ammeter:ammeter}{ऐमीटर} की तरह यह भी
``ताक़तवर लगता है'' वाली एक और ज़मीन सेवामुक्त करके ईमानदार अंकों के
हवाले कर देता है।
\end{definition}

\begin{method}[किसी बल को मापना]\label{met:g9:weight-vs-mass:measure}
\begin{enumerate}
  \item ऐसा \omterm{def:g9:weight-vs-mass:dynamometer}{बलमापी} चुनो जिसका परास काम के मुताबिक़ हो --- क़लमदानों के
        लिए \qty{2}{N} वाला यंत्र, और बस्तों के लिए \qty{100}{N}
        वाला;
  \item उसे उसी हालत में शून्य पर लाओ जिसमें उसे बरतना है (लटकाकर,
        अगर बोझ लटकने वाला हो --- कमानी का अपना \omterm{def:g4:weight-and-mass:weight}{भार} शून्य खिसका देता
        है);
  \item \omterm{def:g2:pushes-and-pulls:force}{बल} टिककर, यंत्र के \omterm{def:g6:sun-earth-moon:axisorbit}{अक्ष} के साथ-साथ लगाओ, और आँख के बराबर से
        पढ़ो;
  \item \omterm{def:g6:measurement-in-science:measuring}{मात्रक} समेत दर्ज करो --- और चूँकि बलों की दिशाएँ होती हैं,
        जहाँ यह मायने रखे वहाँ दिशा भी दर्ज करो।
\end{enumerate}
चित्रों में मापा हुआ \omterm{def:g2:pushes-and-pulls:force}{बल} एक तीर बन जाता है, जिसे घोषित पैमाने पर खींचा
जाता है --- ``दस \omterm{def:g9:weight-vs-mass:newton}{न्यूटन} प्रति \omterm{def:g2:measuring-length:units}{सेंटीमीटर}'' --- और जो उसी ओर ताना जाता
है जिस ओर \omterm{def:g2:pushes-and-pulls:force}{बल}: यानी बचपन वाले बल-तीर, अब मात्रात्मक होकर।
\end{method}

\section{भार का सूत्र}

\begin{proposition}[भार द्रव्यमान के समानुपाती है]\label{prop:g9:weight-vs-mass:formula}
किसी भी दी हुई जगह पर वस्तु का \omterm{def:g4:weight-and-mass:weight}{भार} $P$ (\unit{N} में) उसके \omterm{def:g3:measuring-mass:mass}{द्रव्यमान}
$m$ (\unit{kg} में) के समानुपाती होता है:
\[
P = m \times g ,
\]
जहाँ $g$, यानी उस जगह की \emph{गुरुत्वीय प्रबलता}\index{गुरुत्वीय
प्रबलता}, बताती है कि हर \omterm{def:g3:measuring-mass:units}{किलोग्राम} सामान पर कितने \omterm{def:g9:weight-vs-mass:newton}{न्यूटन} का खिंचाव
पड़ता है। पृथ्वी की सतह के पास,
\[
g \approx \qty{9.8}{N/kg}:
\]
यानी हर \omterm{def:g3:measuring-mass:units}{किलोग्राम} दस से ज़रा कम \omterm{def:g9:weight-vs-mass:newton}{न्यूटन} से खींचा जाता है। यह सूत्र
पिछले अध्याय का नियम ही है, स्थानीय पोशाक में: $g$ पृथ्वी के \omterm{def:g3:measuring-mass:mass}{द्रव्यमान}
और उसकी त्रिज्या को बाँधकर ``यहाँ'' के लिए एक अंक बना देती है।
\end{proposition}

\begin{omfigure}
\begin{tikzpicture}
\begin{axis}[omaxis, width=10.2cm, height=5.8cm,
    xlabel={द्रव्यमान $m$ (\unit{kg})},
    ylabel={भार $P$ (\unit{N})},
    xtick={0,1,2,3,4,5}, ytick={0,9.8,19.6,29.4,39.2,49.0},
    yticklabels={$0$,$9.8$,$19.6$,$29.4$,$39.2$,$49.0$},
    ymin=0, ymax=54, xmin=0, xmax=5.5]
  \addplot[very thick, omDef] coordinates {(0,0) (5.2,50.96)};
  \addplot[only marks, mark=*, omThm] coordinates
    {(1,9.8) (2,19.6) (3,29.4) (5,49.0)};
  \node[right, font=\small, omDef] at (axis cs:2.6,20.4)
    {ढाल: $g = \qty{9.8}{N/kg}$};
\end{axis}
\end{tikzpicture}

{\small \omterm{def:g9:weight-vs-mass:dynamometer}{बलमापी} और अंकित द्रव्यमानों के एक सेट से मापा गया, \omterm{def:g3:measuring-mass:mass}{द्रव्यमान}
के सामने \omterm{def:g4:weight-and-mass:weight}{भार}: मूल बिंदु से गुज़रती एक सीधी रेखा, जिसका खड़ापन ही उस जगह
की $g$ \emph{है}।}
\end{omfigure}

\begin{example}[सूत्र काम पर]\label{ex:g9:weight-vs-mass:atwork}
\qty{60}{kg} का विद्यार्थी: $P = 60 \times 9.8 = \qty{588}{N}$.
\qty{450}{g} की फ़ुटबॉल: $P = 0.45 \times 9.8 \approx \qty{4.4}{N}$
--- \omterm{def:g3:measuring-mass:units}{किलोग्राम} का ध्यान रखो। उलटा चलाओ: \qty{343}{N} \omterm{def:g4:weight-and-mass:weight}{भार} वाले संदूक का
\omterm{def:g3:measuring-mass:mass}{द्रव्यमान} $m = P/g = 343 \div 9.8 = \qty{35}{kg}$. और उलटे से भी उलटा:
किसी पहाड़ पर जहाँ \qty{10}{kg} के थैले का \omterm{def:g4:weight-and-mass:weight}{भार} \qty{97.7}{N} निकले,
वहाँ की स्थानीय $g$ $97.7 \div 10 = \qty{9.77}{N/kg}$ है --- ऊँचाई के
साथ $g$ ज़रा-सी पतली पड़ जाती है, ठीक जैसा प्रतिलोम वर्ग ने वादा किया
था।
\end{example}

\begin{example}[एक सूटकेस, कई भार]\label{ex:g9:weight-vs-mass:planets}
हर दुनिया अपनी $g$ ख़ुद तय करती है, और अंतरिक्ष-यात्रियों वाली पुरानी
कहानी एक सारणी बन जाती है। तुम्हारा वफ़ादार \qty{10}{kg} का सूटकेस:
\begin{center}
\begin{tabular}{l|c|c}
जगह & $g$ (\unit{N/kg}) & सूटकेस का \omterm{def:g4:weight-and-mass:weight}{भार} \\
\hline
पृथ्वी & $9.8$ & \qty{98}{N} \\
\omterm{def:g2:moon-in-the-sky:moon}{चंद्रमा} & $1.6$ & \qty{16}{N} \\
मंगल & $3.7$ & \qty{37}{N} \\
बृहस्पति की बादल-चोटियाँ & $24.8$ & \qty{248}{N}
\end{tabular}
\end{center}
हर जगह दस \omterm{def:g3:measuring-mass:units}{किलोग्राम} सामान --- और हर दुनिया पर अलग \omterm{def:g4:weight-and-mass:weight}{भार}: \omterm{def:g3:measuring-mass:mass}{द्रव्यमान} जस का
तस रहता है, जबकि \omterm{def:g4:weight-and-mass:weight}{भार} स्थानीय $g$ के पीछे चलता है।
\end{example}

\begin{remark}[दोनों तराज़ुओं का फ़ैसला, हमेशा के लिए]\label{rem:g9:weight-vs-mass:scales}
अब बचपन के फ़ैसले सूत्रों के साथ बंद होते हैं। \omterm{def:g3:levers-and-scales:scale}{तुला} $m_1 g$ की तुलना
$m_2 g$ से करती है: स्थानीय $g$ दोनों पलड़ों को गुणा करती है और कट जाती
है --- इसलिए तुलाएँ \emph{\omterm{def:g3:measuring-mass:mass}{द्रव्यमान}} मापती हैं, हर दुनिया पर सच्ची।
जबकि कमानी वाली तराज़ू सीधे $P = m g$ पढ़ती है: वह \emph{\omterm{def:g4:weight-and-mass:weight}{भार}} मापती है,
और अपनी जगह की वफ़ादार रहती है। स्नानघर वाली तराज़ू भेस बदले कमानी वाली
तराज़ू ही है, जिसे कारख़ाने में तुम्हारे \omterm{def:g4:weight-and-mass:weight}{भार} को \emph{पृथ्वी} की $g$ से
भाग देकर \omterm{def:g3:measuring-mass:units}{किलोग्राम} में अंकित कर दिया गया है --- घर पर ईमानदार, \omterm{def:g2:moon-in-the-sky:moon}{चंद्रमा}
पर चापलूस (तुम्हारे हर \qty{60}{kg} के बदले \qty{10}{kg} पढ़ती हुई),
और बृहस्पति पर बदनाम करने वाली।
\end{remark}

\section{काग़ज़ पर बल}

\begin{example}[मेज़ पर रखी किताब, न्यूटनों में]\label{ex:g9:weight-vs-mass:book}
इस पुस्तक का सबसे पुराना संतुलन, आख़िरकार अंकों के साथ: \qty{2}{kg} का
शब्दकोश एक मेज़ पर टिका है। उसका \omterm{def:g4:weight-and-mass:weight}{भार}: $P = 2 \times 9.8 =
\qty{19.6}{N}$, जिसे दस \omterm{def:g9:weight-vs-mass:newton}{न्यूटन} प्रति \omterm{def:g2:measuring-length:units}{सेंटीमीटर} के पैमाने पर
\qty{2}{cm} लंबे तीर की तरह खींचा जाता है, और जो किताब से सीधे नीचे
ताना होता है। मेज़ का सँभालने वाला धक्का: सीधे ऊपर \qty{19.6}{N} ---
यानी बराबर तीर, उलटी दिशा में। बराबरी ठीक-ठीक है, और ठहराव उसका सबूत:
बचपन वाला चित्र, अब मापा हुआ बहीखाता।
\end{example}

\begin{remark}[बल करते क्या हैं --- वह छूटा हुआ नियम]\label{rem:g9:weight-vs-mass:dynamics}
यह अध्याय \omterm{def:g2:pushes-and-pulls:force}{बल} मापता है; यह अभी यह नहीं कहता कि \omterm{def:g2:pushes-and-pulls:force}{बल} का \emph{अधिशेष} गति
के साथ क्या करता है। बचपन के चारों काम --- चलाना, रोकना, मोड़ना, आकार
बदलना --- अपने ठीक-ठीक नियम का इंतज़ार कर रहे हैं, और वह पूरी भौतिकी के
सबसे महान नियमों में से एक है: हाई स्कूल वाली जिल्द अपनी यांत्रिकी
उसी से खोलती है। यह साल बुनियाद पूरी करता है: न्यूटनों में \omterm{def:g2:pushes-and-pulls:force}{बल}, सूत्र से
\omterm{def:g4:weight-and-mass:weight}{भार}, और पैमाने पर खींचे गए तीर --- यानी आने वाले उस वाक्य का व्याकरण।
\end{remark}

\section{अभ्यास}

\begin{exercise}[$\star$]\label{exo:g9:weight-vs-mass:1}
\omterm{def:g2:pushes-and-pulls:force}{बल} का \omterm{def:g6:measurement-in-science:measuring}{मात्रक} और उसका यंत्र बताओ। एक \omterm{def:g9:weight-vs-mass:newton}{न्यूटन} के लिए और सौ \omterm{def:g9:weight-vs-mass:newton}{न्यूटन} के लिए
दो-दो रोज़मर्रा लंगर दो।
\end{exercise}

\begin{exercise}[$\star$]\label{exo:g9:weight-vs-mass:2}
\omterm{def:g4:weight-and-mass:weight}{भार} का सूत्र मात्रकों समेत लिखो, और $g$ का मतलब भी शब्दों में।
\end{exercise}

\begin{exercise}[$\star$]\label{exo:g9:weight-vs-mass:3}
पृथ्वी पर \omterm{def:g4:weight-and-mass:weight}{भार} निकालो: \qty{25}{kg} का कुत्ता; \qty{700}{g} की रोटी;
और ख़ुद तुम (अपना \omterm{def:g3:measuring-mass:mass}{द्रव्यमान} चुन लो)।
\end{exercise}

\begin{exercise}[$\star$]\label{exo:g9:weight-vs-mass:4}
लटकता खरबूज़ा \omterm{def:g9:weight-vs-mass:dynamometer}{बलमापी} को \qty{14.7}{N} तक खींच देता है। उसका \omterm{def:g3:measuring-mass:mass}{द्रव्यमान}?
\end{exercise}

\begin{exercise}[$\star$]\label{exo:g9:weight-vs-mass:5}
अध्याय का आलेख पढ़ो: उसकी ढाल का क्या मतलब है, और \omterm{def:g2:moon-in-the-sky:moon}{चंद्रमा} पर वह रेखा
क्या करती?
\end{exercise}

\begin{exercise}[$\star$]\label{exo:g9:weight-vs-mass:6}
\qty{10}{kg} का सूटकेस पृथ्वी, \omterm{def:g2:moon-in-the-sky:moon}{चंद्रमा}, मंगल और बृहस्पति की सैर करता
है। सारणी के कौन-से अंक कभी नहीं बदलते, और कौन-से बदलते हैं --- और
क्यों?
\end{exercise}

\begin{exercise}[$\star$]\label{exo:g9:weight-vs-mass:7}
\omterm{def:g3:levers-and-scales:scale}{तुला} हर दुनिया पर सच क्यों बोलती है जबकि स्नानघर वाली तराज़ू \omterm{def:g2:moon-in-the-sky:moon}{चंद्रमा} पर
चापलूसी करती है? सूत्रों से बताओ, एहसासों से नहीं।
\end{exercise}

\begin{exercise}[$\star$]\label{exo:g9:weight-vs-mass:8}
\qty{2}{kg} वाले शब्दकोश के दोनों बल-तीर पाँच \omterm{def:g9:weight-vs-mass:newton}{न्यूटन} प्रति \omterm{def:g2:measuring-length:units}{सेंटीमीटर}
के पैमाने पर खींचो (या बयान करो): लंबाइयाँ, दिशाएँ, और बहीखाते का
फ़ैसला।
\end{exercise}

\begin{exercise}[$\star\star$]\label{exo:g9:weight-vs-mass:9}
कोई हवाई कंपनी हर थैले पर \qty{23}{kg} की इजाज़त देती है। एक यात्री
बहस करता है: ``\omterm{def:g2:moon-in-the-sky:moon}{चंद्रमा} पर तो इसका \omterm{def:g4:weight-and-mass:weight}{भार} कुछ ही \omterm{def:g9:weight-vs-mass:newton}{न्यूटन} होता!'' ईंधन और
फ़र्श की मज़बूती को छोड़ भी दें, तो कंपनी की सीमा का सरोकार ठीक-ठीक
\omterm{def:g3:measuring-mass:mass}{द्रव्यमान} से ही क्यों है, और \omterm{def:g2:moon-in-the-sky:moon}{चंद्रमा} के हवाई अड्डे पर उसे कौन-सा यंत्र
जाँचना चाहिए?
\end{exercise}

\begin{exercise}[$\star\star$]\label{exo:g9:weight-vs-mass:10}
मंगल पर किसी रोवर की बाँह का \omterm{def:g4:weight-and-mass:weight}{भार} \qty{111}{N} है। उसका \omterm{def:g3:measuring-mass:mass}{द्रव्यमान}? और
वही बाँह वापस पृथ्वी के जोड़-कक्ष में कितने \omterm{def:g4:weight-and-mass:weight}{भार} की होगी?
\end{exercise}

\begin{exercise}[$\star\star$]\label{exo:g9:weight-vs-mass:11}
सपाट लिटाकर शून्य किया गया \omterm{def:g9:weight-vs-mass:dynamometer}{बलमापी}, फिर लटकाकर बरता जाए, तो थोड़ा झूठ
बोलता है --- किस ओर, और क्यों? (बोझ के साथ कौन-सा फ़ालतू \omterm{def:g4:weight-and-mass:weight}{भार} जुड़
जाता है?)
\end{exercise}

\begin{exercise}[$\star\star\star$]\label{exo:g9:weight-vs-mass:12}
वह प्रयोग डिज़ाइन करो जो अध्याय का आलेख खींचकर $g$ निकाल दे: सामान,
मापनों की सारणी, आलेख, और ढाल पढ़ने का तरीक़ा। फिर समझाओ कि वही तरीक़ा
\omterm{def:g2:moon-in-the-sky:moon}{चंद्रमा} की किसी प्रयोगशाला में चलाने पर क्या देगा --- कौन-से अंक
बदलेंगे, और कौन-सी विधि अछूती बची रहेगी।
\end{exercise}

\section{समस्या: अंतरग्रहीय साज़-सामान की दुकान}

\begin{problem}[{सप्ताहांत समस्या --- अंतरग्रहीय साज़-सामान की दुकान का
नाप-कक्ष; एक शरीर, चार दुनियाएँ; तराज़ुएँ, बलमापी और सामान की
मेज़}]\label{pb:g9:weight-vs-mass:1}
यह दुकान \omterm{def:g4:solar-system:family}{सौर मंडल} के बंदरगाहों के लिए यात्रियों को लैस करती है
($g$, \unit{N/kg} में: पृथ्वी $9.8$; \omterm{def:g2:moon-in-the-sky:moon}{चंद्रमा} $1.6$; मंगल $3.7$;
बृहस्पति का बादल-पटल $24.8$)। नाप-कक्ष तुम सँभालते हो।

\medskip\noindent\textbf{भाग I --- नाप लेना।}
\qty{70}{kg} \omterm{def:g3:measuring-mass:mass}{द्रव्यमान} वाला एक यात्री सामने आता है।
\begin{enumerate}
  \item चारों बंदरगाहों में से हर एक पर उनका \omterm{def:g4:weight-and-mass:weight}{भार}?
  \item किस बंदरगाह पर वह चिरपरिचित स्नानघर वाली तराज़ू (\omterm{def:g3:measuring-mass:units}{किलोग्राम} में
        अंकित, पृथ्वी के लिए अंशांकित) उनका \omterm{def:g3:measuring-mass:mass}{द्रव्यमान} सही पढ़ती है ---
        और बाक़ी तीनों पर वह क्या पढ़ती है? (हर \omterm{def:g4:weight-and-mass:weight}{भार} को पृथ्वी की $g$
        से भाग दो।)
  \item दुकान की ईमानदार तराज़ू अंकित द्रव्यमानों वाली \omterm{def:g3:levers-and-scales:scale}{तुला} है। वह हर
        बंदरगाह पर क्या बताती है, और क्यों?
  \item यात्री के जूते ज़मीन पर \qty{2000}{N} से ज़्यादा दबाव कभी न
        डालें (केंद्रों के फ़र्श पतले होते हैं)। क्या कोई बंदरगाह ऐसा
        है जहाँ जूते-सहित-शरीर (कुल \qty{70}{kg}) यह नियम तोड़ देता
        है?
\end{enumerate}

\medskip\noindent\textbf{भाग II --- सामान की मेज़।}
\begin{enumerate}[resume]
  \item मंगल जाने वाले संदूक का \omterm{def:g3:measuring-mass:mass}{द्रव्यमान} \qty{40}{kg} है। रवानगी
        (पृथ्वी) पर और पहुँचने (मंगल) पर उसका \omterm{def:g4:weight-and-mass:weight}{भार} निकालो।
  \item क़ुलियों का संघ हर बंदरगाह पर हाथ से ढोए जाने वाले \emph{\omterm{def:g4:weight-and-mass:weight}{भार}}
        की सीमा \qty{200}{N} रखता है। हर बंदरगाह पर कोई क़ुली
        ज़्यादा से ज़्यादा कितना \emph{\omterm{def:g3:measuring-mass:mass}{द्रव्यमान}} ढो सकता है?
  \item ``\omterm{def:g2:moon-in-the-sky:moon}{चंद्रमा} पर \qty{150}{N}'' लिखा एक क्रेट पृथ्वी के सीमा-शुल्क
        से \omterm{def:g3:measuring-mass:mass}{द्रव्यमान} के हिसाब से पास होना है। बदलाव करके बताओ।
  \item मेज़ का \omterm{def:g9:weight-vs-mass:dynamometer}{बलमापी} पृथ्वी पर लटके एक थैले का \omterm{def:g4:weight-and-mass:weight}{भार} \qty{294}{N}
        बताता है। सूची-पत्र के लिए उसका \omterm{def:g3:measuring-mass:mass}{द्रव्यमान}, और \omterm{def:g2:moon-in-the-sky:moon}{चंद्रमा} पर उसका
        \omterm{def:g4:weight-and-mass:weight}{भार}?
\end{enumerate}

\medskip\noindent\textbf{भाग III --- शिकायतें और जिज्ञासाएँ।}
\begin{enumerate}[resume]
  \item एक पहलवान \omterm{def:g2:moon-in-the-sky:moon}{चंद्रमा} से गुस्से में लौटता है: ``छह गुना ताक़त की
        कमाई, रातोंरात ग़ायब!'' सूत्र से उसे तसल्ली दो: हर बंदरगाह पर
        असल में क्या बदला, और क्या कभी नहीं बदला?
  \item बृहस्पति के पटल का रसोइया शिकायत करता है कि फेंटनी चलाना
        निहाई उठाने जैसा लगता है, जबकि फेंटनी का ``\omterm{def:g3:measuring-mass:mass}{द्रव्यमान} महज़
        \qty{200}{g}'' है। पटल पर उसका \omterm{def:g4:weight-and-mass:weight}{भार} निकालो और फ़ैसला दो।
  \item तोहफ़ों की दुकान ``एक \omterm{def:g9:weight-vs-mass:newton}{न्यूटन} चॉकलेट, कहीं भी'' बेचती है। किस
        बंदरगाह पर ख़रीदार को सबसे ज़्यादा चॉकलेट \emph{सामान} मिलता
        है, और पृथ्वी, \omterm{def:g2:moon-in-the-sky:moon}{चंद्रमा} तथा बृहस्पति में से हर एक पर कितना
        \omterm{def:g3:measuring-mass:mass}{द्रव्यमान}?
  \item दुकान का आदर्श वाक्य लिखो: एक वाक्य, जो यह अलग कर दे कि
        तुम्हारे साथ क्या सफ़र करता है और हर दुनिया उस पर क्या चुंगी
        लगाती है।
\end{enumerate}
\end{problem}
